\chapter{Verification and Validation Methodology}
\label{chap:validation-method}

% ============================================================================
% CHAPTER PURPOSE -- VERY HIGH PRIORITY
% Define HOW correctness and credibility will be assessed. Keep results out of
% this chapter except for tiny illustrative checks. Chapter \ref{chap:results}
% reports outcomes.
%
% Distinguish terminology:
% - software verification: was the model implemented correctly?
% - model verification: do equations/logical behaviours behave as intended?
% - validation: does the model adequately represent observed race behaviour for
%   its stated purpose and operating domain?
% ============================================================================

\section{Evaluation Framework}
\label{sec:evaluation-framework}

% Introduce the chain of evidence:
% equation checks -> software tests -> limiting/expected-behaviour tests ->
% historic-case validation -> sensitivity/uncertainty analysis.
% SUGGESTED FIGURE 5.1: verification and validation workflow.
% Map each method to a success criterion from Section \ref{sec:success-criteria}.

\section{Software and Calculation Verification}
\label{sec:calculation-verification}

\subsection{Unit-Test Design}
% Define test families and expected values, not framework trivia:
% - base lap with all penalties zero;
% - tyre age/degradation increments;
% - compound offset;
% - warm-up only on intended laps;
% - cumulative summation;
% - sign convention for gain;
% - initial gap to post stop gap conversion;
% - common pit loss cancellation;
% - invalid/missing inputs.
% Put exhaustive test-case tables/screenshots in Appendix \ref{app:software-tests}.

\subsection{Integration and Regression Testing}
% Verify data ingestion -> model state -> result -> visual marker.
% Fix at least one known scenario/parameter set as a regression case.
% State pass/fail criteria before reporting results.

\subsection{Independent Hand Calculation}
% HIGH VALUE. Recalculate the worked example independently (spreadsheet/manual)
% and compare it with software output. This is stronger evidence than tests that
% reproduce the same implementation logic.

\section{Model Behaviour Verification}
\label{sec:model-behaviour-verification}

% Test limiting and monotonic behaviours expected from the formulation:
% - zero degradation and equal compounds -> near-zero tyre-driven gain;
% - increased target degradation -> stronger undercut, all else equal;
% - larger warm-up penalty -> weaker undercut;
% - traffic penalty -> weaker undercut;
% - zero offset laps -> no offset-lap gain;
% - identical scenarios -> zero gain;
% - used replacement tyre behaves consistently with starting age.
% These tests answer "does the model make engineering sense?" before historic
% validation.

\section{Historic Validation Case Selection}
\label{sec:case-selection}

% VERY HIGH PRIORITY. Define transparent selection criteria to avoid cherry
% picking. Aim for a small but varied set rather than one famous success.
% Candidate categories:
% - clear successful undercut;
% - unsuccessful/neutral attempt;
% - case compromised by traffic/warm-up;
% - different circuits/compounds/offset lengths;
% - optional negative-control period with no expected opportunity.
% Explain data completeness and why the cases are observable.
% Record excluded cases and reasons in Appendix \ref{app:validation-data}.

\section{Historic Data Preparation}
\label{sec:validation-data-prep}

% Define consistently:
% - decision lap and common comparison endpoint;
% - actual tyre ages/compounds;
% - treatment of in/out laps;
% - removal/retention of laps affected by flags, traffic or errors;
% - base pace and degradation estimation method;
% - parameter values unavailable directly from data;
% - actual observed gain metric.
% This section is crucial for reproducibility and prevents hindsight bias.

\section{Validation Metrics}
\label{sec:validation-metrics}

% Define metrics before seeing/reporting conclusions. Possible metrics:
% - predicted and observed gain in seconds;
% - signed error e = predicted - observed;
% - absolute error;
% - correct direction/classification of gain/loss;
% - post stop gap error;
% - opportunity-lap timing error if detecting a window.
% For a small sample, do not overstate statistical power. Report each case and
% summary statistics transparently.

\section{Sensitivity Analysis Method}
\label{sec:sensitivity-method}

% VERY HIGH PRIORITY / likely major mark source.
% Establish:
% - baseline case/parameter set;
% - parameter ranges and justification;
% - one-at-a-time (OAT) method for interpretability;
% - optional multi-parameter scenarios/interactions if time permits;
% - output metric (usually change in G or threshold/crossover).
% Minimum parameters:
% - degradation rate/current tyre age;
% - compound offset;
% - replacement tyre warm-up;
% - number of offset laps;
% - relative driver/car base pace;
% - traffic penalty if retained.
% Pit loss should be treated with the symmetry discussion from Section
% \ref{sec:pit-loss}; test unequal/rejoin effects only if relevant.
%
% MEDIUM/LOW PRIORITY extension:
% factorial/global sensitivity or Monte Carlo uncertainty. Do not add this at
% the expense of a complete, well-explained OAT analysis.

\section{Uncertainty and Error Treatment}
\label{sec:uncertainty-method}

% HIGH VALUE. Identify uncertainty sources:
% - noisy/traffic-affected historic lap times;
% - degradation fit uncertainty;
% - unknown warm-up/traffic penalty;
% - driver/car pace changes;
% - API missing/incorrect fields.
% State whether uncertainty is represented by ranges, alternate assumptions,
% confidence intervals, or qualitative caveats.
% Avoid claiming precise predictive certainty from manually selected inputs.

\section{Validation Threats and Bias Control}
\label{sec:validation-threats}

% MEDIUM/HIGH VALUE.
% Address:
% - hindsight in selecting known undercuts;
% - circularity if parameters are fitted using the same outcome being predicted;
% - confounding from driver pace, fuel, traffic and track evolution;
% - small case count/generalisation;
% - data-source limitations.
% Explain mitigations such as fixed selection criteria, holdout cases, baseline
% extraction before examining the outcome, and reporting failed cases.

\section{Software Usability Check}
\label{sec:usability-check}

% LOW PRIORITY. Only include if actually done.
% A short task-based check can confirm that a user can load a race, configure a
% scenario and interpret an annotation. Do not imply formal human-subject UX
% research without appropriate design/ethics.

\section{Chapter Summary}
\label{sec:validation-method-summary}

% Summarise the evidence that will be reported in Chapter \ref{chap:results}.
